% method_sim.tex
% Method: variance-corrected SIM composition.
% Simplified 2026-07-14 around a direct variance explanation. The
% section now proceeds from the composition goal, to naive double
% counting, to the variance correction, the high-R^2 exception, and
% finally the evaluation protocol. Detailed verification remains in
% the Online Appendix.

The single-asset model generated each ticker separately and did not
specify cross-asset dependence. For the multi-asset model, we added a
shared market path as a common factor, using the single-index model
(SIM) of Sharpe~\cite{sharpe1963}. We write $g_i(t)$ for the growth rate $G_{i,t}$ of asset $i$ on days
$t=1,\dots,T$. The horizon $T$ is $2{,}766$ days in the training
comparisons and $249$ days in the holdout. The same symbols denote
observed paths when we estimate parameters and synthetic paths when
we compose them:
\begin{equation}
  g_i(t) \;=\; \alpha_i \;+\; \beta_i\, g_m(t) \;+\; \eps_i(t),
  \qquad t = 1,\dots,T.
  \label{eq:sim-composition}
\end{equation}
Here $\alpha_i$ is the intercept for asset $i$, $g_m$ is the
growth-rate path of the market, SPY, $\beta_i$ is the asset's market
loading, and $\eps_i$ is its asset-specific residual path. We use
$\eps_i$ for the residual path and $\epsilon$ for the jump probability
in the single-asset model. We estimated the intercept and market
loading $(\alpha_i,\beta_i)$ from observed data by ordinary least
squares (OLS). The choice of $\eps_i$ determined the composition
method. A standard SIM draws independent Gaussian residuals with the
variance estimated by OLS. These Gaussian residuals preserve the fitted market relationship but
lack the heavy tails, and with jumps enabled the volatility
clustering, that the single-asset generator produces. We called the Gaussian-residual method
\emph{Gaussian SIM} and used it as one of the compared methods.

\paragraph{Naive full-path composition.}

Using a full synthetic growth-rate path $\tg_i$ drawn from the asset's
own single-asset generator, the \emph{generator draw}, directly as the
residual would give:
\begin{equation}\label{eq:naive-composition}
g_i = \alpha_i + \beta_i\, g_m + \tg_i.
\end{equation}
The fitted intercept and market term already set the conditional mean
of asset $i$, while the generator draw generally has
$\E[\tg_i]\ne 0$. Equation~\eqref{eq:naive-composition} therefore has
conditional mean $\alpha_i+\E[\tg_i]+\beta_i g_m$, not
$\alpha_i+\beta_i g_m$.
We therefore centered every generator draw before composition,
subtracting its realized path mean:
\begin{equation}
  \overline{\tg}_i \equiv \frac{1}{T}\sum_{t=1}^{T}\tg_i(t),
  \qquad
  \tg_i^{\mathrm c}(t) \equiv \tg_i(t)-\overline{\tg}_i.
  \label{eq:center-generator}
\end{equation}
Therefore, $T^{-1}\sum_t\tg_i^{\mathrm c}(t)=0$ exactly. Subtracting
the realized rather than the population mean makes the centering exact on every path. Subtracting a
constant leaves both the sample variance and the sample covariance with
$g_m$ unchanged. We used the centered construction
$g_i=\alpha_i+\beta_i g_m+\tg_i^{\mathrm c}$, which corrects the mean
but not the variance, as the \emph{naive} benchmark in all
comparisons. The sample variance of each draw,
$\sigma_{\gen,i}^2=\Var(\tg_i)=\Var(\tg_i^{\mathrm c})$, is the
\emph{generator variance}; it is recomputed for every draw. It already represents the asset's total
variance, so adding the market term adds market variance a second
time. During composition, $g_m$ was the observed SPY path in the training
comparison and a path simulated from SPY's closing-price generator in the holdout
and jump comparisons, with one market path per replication shared by
all assets, and
in both cases it was independent of the asset draws. The covariance
of a centered draw with the market path is zero in population and close to zero on a long realized
path. We approximated the naive composed variance accordingly:
\begin{equation}
  \Var(g_i)
  \;\approx\; \beta_i^2\sigma_m^2 + \sigma_{\gen,i}^2,
  \qquad \Cov(\tg_i^{\mathrm c},g_m)\approx 0,
  \label{eq:naive-inflation}
\end{equation}
where $\sigma_m^2=\Var(g_m)$ is the sample variance of the market path
used in composition. The naive composed variance therefore exceeds the generator variance
by $\beta_i^2\sigma_m^2$, an excess that grows with the square of the
market loading.

\paragraph{Variance correction.}

We set $\sigma_{\gen,i}^2$ as the target total variance for asset $i$.
The market term contributed $\beta_i^2\sigma_m^2$, so the residual term
could contribute only the remaining variance. We defined $s_i$ as the
nonnegative asset-specific residual scale, set
$\eps_i=s_i\tg_i^{\mathrm c}$, and required
$\Var(\beta_i g_m+s_i\tg_i^{\mathrm c})=\sigma_{\gen,i}^2$. Under the
zero-covariance assumption, solving for $s_i$ gives the correction:
\begin{equation}
  \boxed{\;s_i^2 \;=\; 1 \;-\; \frac{\beta_i^2\, \sigma_m^2}{\sigma_{\gen,i}^2}\;}
  \label{eq:scale}
\end{equation}
We call the market term's share of the target variance the
\emph{market loading ratio}:
\begin{equation}
  \rho_i \;\equiv\; \frac{\beta_i^2\, \sigma_m^2}{\sigma_{\gen,i}^2},
  \label{eq:rho}
\end{equation}
so that $s_i^2=1-\rho_i$ under the zero-covariance assumption. At
$\rho_i=1$, the rule of Eq.~\eqref{eq:scale} leaves no asset-specific
variance; for $\rho_i>1$, it has no real-valued solution.

With the observed variances, $\rho_i$ would equal the coefficient of
determination of the regression, which is below one. Because $\rho_i$
instead uses the generator and market-path variances, it can reach or exceed one when the generator
variance is low or the market path is more volatile than the observed
SPY training series, particularly for assets with large $\beta_i$. We therefore required the asset-specific draw to
contribute at least a fraction $f$ of the target variance, the
\emph{floor}, and set
$\rho_i^{\max}=1-f$. We applied the correction through an
\emph{effective loading} $\beta_i^{\eff}$, the loading used in
composition, which equals $\beta_i$ unless $\rho_i>\rho_i^{\max}$:
\begin{equation}
  \beta_i^{\eff} \;=\;
  \begin{cases}
    \beta_i & \text{if } \rho_i \le \rho_i^{\max} \\[4pt]
    \sign(\beta_i)\,\sqrt{(1 - f)\,\dfrac{\sigma_{\gen,i}^2}{\sigma_m^2}} & \text{if } \rho_i > \rho_i^{\max}
  \end{cases}
  \qquad
  s_i^2 \;=\;
  \begin{cases}
    1 - \rho_i & \text{if } \rho_i \le \rho_i^{\max} \\[4pt]
    f          & \text{if } \rho_i > \rho_i^{\max}
  \end{cases}
  \label{eq:clip}
\end{equation}
We call the second case of Eq.~\eqref{eq:clip} the \emph{clipped}
case. In it, the correction reduced the magnitude of the loading but
kept the sign of the calibrated $\beta_i$. We used $f=0.10$, so the asset-specific draw contributed
at least $10\%$ of the target variance for every asset under this rule
(Table~\ref{tab:hyperparams}). We retained the fitted intercept
$\alpha_i$ when clipping changed the loading. Because the scaled
residual had zero sample mean, the composed mean was
$\overline{g}_i=\alpha_i+\beta_i^{\eff}\overline{g}_m$, where
$\overline{g}_m$ is the mean of the market path. Clipping therefore
changed the mean by $(\beta_i^{\eff}-\beta_i)\overline{g}_m$ relative
to the calibrated linear mean on that same market path. The clipped case thus kept the intercept and the target variance and
set the residual share to exactly $f$, but changed the loading and,
unless $\overline{g}_m=0$, the mean. Online
Appendix~\ref{sec:supp-residual-choice} examines the substitution of
a full-return draw for the residual: it derives what the construction
preserves exactly and how much of the generator's excess kurtosis
survives composition as $\rho_i$ grows.

\paragraph{High-$R^2$ index trackers.}

The coefficient of determination $R^2_{i,\real}$, from the OLS
regression of the observed asset on observed SPY, measures the
fraction of the asset's variance explained by the market. The rule of
Eq.~\eqref{eq:clip} targets total generator variance, not
$R^2_{i,\real}$. Under the zero-covariance
assumption, the composed $R^2$ under that rule equals $\rho_i$, capped
at $1-f$ by the floor. It therefore cannot exceed $1-f$, and it matches $R^2_{i,\real}$ only
when $\sigma_m^2/\sigma_{\gen,i}^2$ equals the ratio of the observed SPY
variance to the observed asset variance. Index-tracking funds, such as the Nasdaq-100 fund QQQ and the
S\&P~500 growth fund SPYG, have $R^2$ against SPY close to or above
the cap $1-f$, so that rule could not reliably reproduce their co-movement with the
market. We therefore used two branches: the rule of
Eq.~\eqref{eq:clip}, the \emph{variance-preserving branch}, for most
assets, and a separate \emph{$R^2$-preserving branch} when
$R^2_{i,\real}$ was at least $R^2_{\mathrm{preserve}}$. We set $R^2_{\mathrm{preserve}}=0.80$,
which lies in a wide gap of the observed $R^2$ distribution between
the index trackers and all other assets in our universe
(Fig.~\ref{fig:r2_distribution}). This branch keeps $\beta_i^{\eff}=\beta_i$ and uses the same
market-path variance $\sigma_m^2$ as the variance-preserving branch.
We computed the residual variance required to recover $R^2_{i,\real}$
under the zero-covariance
assumption:
\begin{equation}
  \boxed{\;\sigma_{\eps,\mathrm{target}}^2 \;=\; \beta_i^2\, \sigma_m^2 \cdot
    \frac{1 - R^2_{i,\real}}{R^2_{i,\real}}\;}
  \label{eq:r2-target}
\end{equation}
We rescaled the generator draw to the target residual variance:
\begin{equation}
  \eps_i \;=\; \left(\sqrt{\frac{\sigma_{\eps,\mathrm{target}}^2}{\sigma_{\gen,i}^2}}\,\right)\tg_i^{\mathrm c}.
  \label{eq:r2-rescale}
\end{equation}
This branch preserves $R^2_{i,\real}$ rather than the generator
variance: the composed variance becomes
$\beta_i^2\sigma_m^2/R^2_{i,\real}$. The scale factor in Eq.~\eqref{eq:r2-rescale} exceeds one when the generator variance is below the target residual
variance; the composed variance then exceeds the generator variance.

\paragraph{Complete composition rule.}

The complete rule has three outcomes, recorded by the branch flag
$b_i$ in Algorithm~\ref{alg:hybrid}. Index trackers used the
$R^2$-preserving branch; all other assets used the
variance-preserving branch, with the calibrated loading or, in the
clipped case, a reduced loading. Each composed path is then
$g_i=\alpha_i+\beta_i^{\eff}g_m+\eps_i$. We call this complete rule
the \emph{hybrid} composition, the label used in the tables and
figures, and all results reported for the variance correction refer
to it. Membership in the $R^2$-preserving branch depends only on
$R^2_{i,\real}$ and is fixed for each asset. Everything else is
computed per draw: because $\sigma_{\gen,i}^2$ is recomputed for each
draw, so are $s_i$, $\rho_i$, and $\beta_i^{\eff}$, and one asset can be
clipped in some draws and not in others. Because the draws are independent of the market path, OLS estimates
from composed paths converge, as the horizon grows, to the intercept
$\alpha_i$ and the effective loading $\beta_i^{\eff}$, which equals
$\beta_i$ unless the draw is clipped. The variance target and the loading recovery are exact only when the
sample covariance between the draw and the market path is zero; on a
finite path, they deviate by that covariance term (Online
Appendix~\ref{sec:supp-sim-derivations}). Because centering and
scaling fix the residual sum, every branch imposes a zero residual sum over the horizon $T$, which is stronger than a zero population
mean. For a fixed market path and effective loading, the terminal
cumulative log return therefore has no asset-specific residual
uncertainty, although daily and intermediate returns still vary. In
the clipped case, the effective loading depends on the variances of
each draw, so it can vary between draws and produce terminal
variation through the market term. Because centering and scaling use
statistics of the entire draw, the construction also depends on the
chosen horizon. We therefore evaluated it as a constrained
finite-horizon construction for daily distributional fit and recovery
of the market loading; it does not preserve cumulative residual uncertainty.

The construction links assets only through same-day market returns.
A large market move can therefore affect several assets together, but
the rule does not separately fit lead--lag effects or a common
volatility factor in the residuals. Algorithm~\ref{alg:hybrid} also
allows an optional rank reorder that adds residual dependence across
assets; the six-method comparison did not use it. The reorder draws
one joint vector per day from a copula fitted to the observed OLS
residuals and permutes each asset's residual values over days so that
their ranks over time match the ranks of that asset's copula
coordinate. The reorder imposes cross-asset residual
dependence and keeps each path's values and variance, but it replaces the generator's day-to-day ordering and with it any
volatility clustering in the draw.

\begin{algorithm}[t]
\caption{Hybrid single-index model (SIM) composition with variance correction.}
\label{alg:hybrid}
\small
\begin{algorithmic}[1]
  \REQUIRE Nonconstant $g_m,\tg_i\in\R^T$, $i=1,\dots,N$; each
    $\tg_i$ is independent of $g_m$.
  \REQUIRE Parameters $(\alpha_i,\beta_i,R^2_{i,\real})$; floor
    $f\in(0,1)$; threshold $R^2_{\mathrm{preserve}}\in(0,1]$.
  \ENSURE Composed paths $g_i\in\R^T$ and branch flags $b_i$.
  \STATE $\sigma_m^2 \gets \Var(g_m)$.
  \FOR{each asset $i = 1,\dots,N$}
    \STATE $\sigma_{\gen,i}^2 \gets \Var(\tg_i)$;
      $\tg_i^{\mathrm c} \gets \tg_i-\overline{\tg}_i$.
    \IF{$R^2_{i,\real} \ge R^2_{\mathrm{preserve}}$}
      \STATE $v_i \gets \beta_i^2\sigma_m^2(1-R^2_{i,\real})/R^2_{i,\real}$.
      \STATE $s_i \gets \sqrt{v_i/\sigma_{\gen,i}^2}$;
        $\beta_i^{\eff} \gets \beta_i$; $b_i\gets\textsc{r2-preserve}$.
    \ELSE
      \STATE $\rho_i \gets \beta_i^2 \sigma_m^2 / \sigma_{\gen,i}^2$
      \IF{$\rho_i \le 1 - f$}
        \STATE $s_i \gets \sqrt{1-\rho_i}$;
          $\beta_i^{\eff}\gets\beta_i$; $b_i\gets\textsc{variance}$.
      \ELSE
        \STATE $s_i\gets\sqrt{f}$;
          $\beta_i^{\eff}\gets\sign(\beta_i)
          \sqrt{(1-f)\sigma_{\gen,i}^2/\sigma_m^2}$.
        \STATE $b_i\gets\textsc{clipped}$.
      \ENDIF
    \ENDIF
    \STATE $\eps_i \gets s_i\tg_i^{\mathrm c}$.
  \ENDFOR
  \STATE \textit{Optional:} jointly rank-reorder $\{\eps_i\}_{i=1}^{N}$.
  \FOR{each asset $i = 1,\dots,N$}
    \STATE $g_i \gets \alpha_i + \beta_i^{\eff}\, g_m + \eps_i$
      \hfill \COMMENT{Eq.~\eqref{eq:sim-composition}}
  \ENDFOR
\end{algorithmic}
\end{algorithm}

% --- Table: Hyperparameters ---
% Main table tab:hyperparams is collected after the references.
\FloatBarrier

\paragraph{Evaluation protocol.}

We evaluated the method on $424$ United States equities and exchange-traded funds
from Polygon.io. Each ticker had $2{,}767$ daily closing prices from
January 2014 through December $2024$, which yielded $2{,}766$ growth rates; we kept only tickers with all
$2{,}767$ prices. This complete-case selection conditioned the analysis on assets
that survived and retained data coverage throughout the training
window. We used observed SPY as the market path in the training
comparison and composed the other $423$ assets. For each asset, including SPY, we fitted the hidden Markov generator used in the SPY analysis
(state partition, counted transitions, and emissions) to the
closing-price growth rates
(Table~\ref{tab:hyperparams}; $r_f=0$). We regressed each observed
asset on observed SPY to estimate $\alpha_i$, $\beta_i$,
$R^2_{i,\real}$, and residual variance. Across the sample, $\beta_i$
ranged from $0.03$ to $1.79$, and $R^2_{i,\real}$ ranged from $0.001$
to $0.933$ (Figs.~\ref{fig:branch_map} and \ref{fig:r2_distribution}).
We compared six ways to construct the asset-specific term
(Table~\ref{tab:aggregate}), which we call the six-method comparison.
In this comparison, we set the jump probability to zero in every
generator, so each was an HMM-NJ, to measure composition without
per-ticker jump calibration. The naive method used the centered
full-return draw without variance correction. The Gaussian SIM method
used independent Gaussian residuals with variance estimated by OLS.
The hybrid method used Algorithm~\ref{alg:hybrid}. The residual-fit
hidden Markov method, labeled JumpHMM-on-residuals in the tables after JumpHMM, the
software package that implements our generator, fitted a new HMM-NJ to the OLS residual series after converting the
residuals to a synthetic price path by cumulative compounding. The
block-bootstrap method sampled the OLS residual series with the
Politis--Romano stationary bootstrap~\cite{politisRomano1994} and a
mean block length of $50$ days, close to $\sqrt{T}$ for the
$2{,}766$-day training series; the holdout used the same block
length. The GARCH(1,1)-$t$ method
fitted a conditional-variance model to each residual
series~\cite{bollerslev1986}. These four methods added their residual
draws to $\alpha_i+\beta_i g_m$ without centering or rescaling. We accepted a GARCH fit only if the
optimizer converged, the likelihood and coefficients were finite, and
the variance intercept was positive. The ARCH and GARCH coefficients
also had to be nonnegative with a sum below one, and the Student-$t$
degrees of freedom had to exceed two. Finally, the unconditional
variance had to be finite and positive, and a trial simulated path had
to be finite and nonconstant. All $423$ optimizations converged; $30$ fits failed the stationarity
condition that the ARCH and GARCH coefficients sum below one, leaving
$393$ eligible training assets. We rechecked each stored fit against
these criteria before evaluation.

For each ticker, the naive and hybrid methods received the same
generator draw, which isolated the effect of the variance correction;
the other four methods generated their own residual draws. We did not
use the optional cross-sectional rank reorder because the experiment
tested each ticker's marginal construction separately. The naive and
hybrid methods force the residual mean to zero over the horizon,
while the other four methods do not. To isolate the effect of this
difference in mean constraints, we also compared each method's native
path with the same path after subtracting its residual mean over the
full training horizon (Table~\ref{tab:centering_control}). Only the residual mean differed between the two versions of each
method; the fitted models, random draws, and market path were the
same. Naive and hybrid
paths already had zero residual mean. The control used $100$
replications per asset and the same $393$ assets, those whose GARCH fits passed the acceptance
criteria, for all six methods. Online
Appendix~\ref{sec:supp-centering-control} gives the control definition
and separates centering effects from model fitting.

We evaluated each composed series on factor recovery, distributional
fit, and variance preservation. For factor recovery, a new OLS
regression of the composed path on the market path produced the
absolute errors $|\hat\alpha-\alpha_i|$, $|\hat\beta-\beta_i|$, and
$|\hat R^2-R^2_{i,\real}|$
against the calibrated values; lower values indicate closer recovery.
On the variance-preserving branch, the composed $R^2$ approaches each
draw's $\rho_i$, capped at $1-f$, rather than $R^2_{i,\real}$, so for
the hybrid method the $R^2$ error mainly measures this designed shift
rather than estimation error.
Because the reference is the calibrated loading, the loading error of
a clipped asset includes the change made by clipping. For distributional fit, we tested each composed path against the
observed returns of the same asset and reported KS and AD pass rates,
Wasserstein-1 distance, excess kurtosis, and the Hill index of the
absolute-return tail. Higher KS and AD pass rates and lower Wasserstein distance indicate a
closer match; for kurtosis and the Hill index, values closer to the observed asset's indicate a closer tail
match. The six-method comparison did not score volatility clustering;
the jump-enabled comparison did. For variance preservation of the naive and hybrid methods, which share
a generator draw, we used the composed-to-generator variance ratio. A
ratio of one indicates equal variances; the $R^2$-preserving branch
does not target this ratio.
Online Appendix~\ref{sec:supp-validation} gives the metric definitions.
We ran $100$ replications per ticker with seed $1234$, producing $42{,}300$ series for each method except GARCH, which
produced $39{,}300$ series across its $393$ eligible assets. Training
summaries pooled these series across tickers.
We archived all $250{,}800$ training-path scores together with records
of the inputs, fitted models, source code, and configuration. The
aggregate, branch, and loading-quartile tables
(Tables~\ref{tab:aggregate}, \ref{tab:branch}, and
\ref{tab:beta_bucket}) were generated from this archived run, with matching method-level summaries distributed
alongside the path scores.

For the six-method holdout evaluation, we froze all marginal fits,
residual models, and SIM parameters estimated from $2014$--$2024$
(Tables~\ref{tab:oos_composition} and \ref{tab:var}). We evaluated
the methods against the $249$ growth rates observed in $2025$. Of the
$423$ non-market training assets, $416$ had complete holdout
histories; GARCH results covered $386$ because its $30$ non-stationary training fits remained excluded
(Table~\ref{tab:oos_composition}). For each replication, we generated
a $249$-day SPY path from the closing-price SPY model with jumps
disabled and used the same simulated market path in all asset
compositions. In this evaluation, the correction used the observed
training-period market variance in place of the simulated path
variance $\sigma_m^2$ in
Eqs.~\eqref{eq:scale}--\eqref{eq:r2-rescale}; the generator variance
was still computed from each asset draw. The residual scale therefore
budgeted for the training market variance, and variation in the
simulated market variance could move the composed variance away from
the generator target even when the sample cross-covariance was zero.
Observed $2025$ SPY and asset returns were used only for scoring. We
repeated the KS, AD, Wasserstein-1, and kurtosis evaluations with
$100$ paths per ticker and measured variance, for all six methods, as the ratio of synthetic
to observed $2025$ variance. Goodness-of-fit tests have less power on
$249$ observations than on $2{,}766$. We therefore also tested the same $249$-day holdout composed paths
against a randomly selected contiguous $249$-day block of the asset's
observed training returns. This length-matched comparison shows how
the tests behave at $249$ observations without any change specific to
$2025$. We first averaged replication-level
results within ticker and then summarized across tickers so each
asset received equal weight.

To isolate each method's one-day tail thresholds from estimation
noise in a single short path, we used a separate simulation ensemble
for the Value-at-Risk (VaR) evaluation (Table~\ref{tab:var}). For each asset and method, we pooled the daily log returns (growth
rates multiplied by $\Delta t$) from $5{,}000$ separately constructed
$249$-day paths, giving $1{,}245{,}000$ observations for the lower-tail
quantile at each of the $95\%$ and $99\%$ coverage levels. Each of the $5{,}000$ replications used one simulated SPY path, shared
across assets and methods. Every hidden Markov generator started from
its fitted stationary state distribution, and centering, scaling, and market variance followed the holdout
conventions, including the training-period market variance. We fixed one threshold per asset, method, and
coverage level before counting exceedances on the observed $2025$ returns. Each threshold is therefore an
unconditional one-day VaR that does
not depend on the simulated state. Before holdout scoring, we fixed the ensemble size and a set of
smaller pool sizes at which we confirmed that the thresholds had
stabilized. As a simulation check, an independent ensemble of $1{,}000$
paths per asset and method compared the simulated probability of
falling below each threshold with its nominal level (Online
Appendix~\ref{sec:supp-var-ensemble}). Because the pool mixes all $249$ days of each path, each threshold is
a quantile of returns pooled over the days of a $249$-day composed
path, not the VaR of any particular day. Pooling did not remove the construction's horizon
dependence or zero-sum residual constraint. We reported mean
exceedance rates and cross-ticker standard deviations descriptively. We did not run a formal coverage
significance test, because
simulated days, dependent holdout observations, and assets sharing a
market are not independent replications. No $2025$ data were used for
fitting or threshold calibration.


\paragraph{Jump-enabled multi-asset comparison.}

To test whether the temporal benefit of the single-asset jump
mechanism survived composition, we compared naive and hybrid
composition under three jump settings (Table~\ref{tab:jump_ablation}).
Unlike the six-method comparison, which used observed SPY in
training, this comparison simulated the SPY market path in both the
$2{,}766$-day training window and the $249$-day holdout, so a jump
setting could apply to the market generator. Jumps were disabled,
enabled only in the SPY market generator, or enabled in both SPY and
every asset generator. We froze the closing-price hidden Markov generators and SIM
calibrations of the six-method comparison. Enabled models
received the SPY settings of Table~\ref{tab:hyperparams},
$(\epsilon,\lambda)=(10^{-4},90)$, $p_{\rm neg}=0.52$, and
$N_{\rm tail}=5$, without per-asset tuning. These settings were tuned on volume-weighted SPY prices and applied
here to generators fitted to closing prices. At these settings,
most paths contain no episode, especially over the $249$-day holdout,
so the jump effects are carried by a minority of paths. In both windows and for all three jump settings, the correction used
each simulated market path's variance and each asset draw's variance.
This differs from the six-method holdout, which fixed the market
variance at its training value.

For each jump setting, naive and hybrid composition received the same
asset and market draws, and one market path was
shared across all assets in each replication. We generated $250$
paths for each of four seeds ($1234$, $2345$, $3456$, and $4567$),
giving $1{,}000$ replications per asset, setting, and method. We used a path-level version of ACF-MAE, computed on each path and
then averaged across paths (Eq.~\eqref{eq:acf-path-error} in Online
Appendix~\ref{sec:supp-validation}); the single-asset ACF-MAE instead
averages the path ACFs before measuring the error. We used
lags $1$--$25$, the lag window of the tuning objective, with lags
$1$--$60$ as a secondary check, and repeated the marginal and variance
diagnostics. The temporal scores of the single-asset and jump comparisons
therefore have different estimands and cannot be compared directly. We included all paths, not only jump-active ones, in the summaries;
jump-active subsets only described the mechanism
(Table~\ref{tab:jump_ablation_episodes}). For Monte Carlo uncertainty, we averaged each metric across tickers
within each replication, which kept assets together under their
shared market path. Jump settings were paired by replication, sharing
the draws of any component whose jump setting did not change, and the
approximate $95\%$ confidence intervals were the mean of the $1{,}000$
replication-level differences plus or minus $1.96$ standard errors
(Table~\ref{tab:jump_ablation_uncertainty}). These intervals describe simulation
uncertainty conditional on the fitted models and observed histories,
not uncertainty across market regimes.
